\documentclass[11pt,a4paper]{article}

\usepackage[T1]{fontenc}
\usepackage[utf8]{inputenc}
\usepackage[english]{babel}        % -> [french]{babel} for a French report

\usepackage{amsmath,amssymb,mathtools}
\usepackage{graphicx}
\usepackage{booktabs}
\usepackage{caption}
\usepackage{siunitx}
\usepackage[margin=2.5cm]{geometry}
\usepackage{microtype}
\usepackage{xcolor}
\usepackage[hidelinks]{hyperref}
\usepackage{cleveref}

% --- artifacts produced by gradient.jl -------------------------------------
% Figures: savefig(...) PDF into figures/.  Tables: PrettyTables LaTeX backend
% into generated/.  Each renders as a placeholder until the script has run, so
% the file always compiles.  Paths are explicit so the existence test and the
% include agree (\graphicspath does not affect \IfFileExists).
\newcommand{\incfig}[2][0.82\linewidth]{%
  \IfFileExists{#2}%
    {\includegraphics[width=#1]{#2}}%
    {\fbox{\parbox[c][0.20\textheight][c]{#1}{\centering\ttfamily
      missing figure: #2\\[3pt]\normalfont\itshape (run gradient.jl)}}}}
\newcommand{\inctbl}[1]{%
  \IfFileExists{#1}{\input{#1}}{{\ttfamily missing table: #1}}}

% visible, removable placeholder for prose written from the measured numbers
\newcommand{\fillin}[1]{\textcolor{gray}{\textit{[#1]}}}

\title{Gradient computation for the inverse problem:\\
       a comparison of differentiation methods}
\author{Sipan Bareyan}
\date{\today}

\begin{document}
\maketitle

\begin{abstract}
\noindent
We compare methods for computing the gradient $g=\nabla_V L$ of the misfit
\(L(V)=\int_0^{2\pi}\bigl(u(V)-u_0\bigr)^2\,dx\), where the state $u(V)$ solves
the discretized nonlinear equation $F(u,V)=0$ by Newton's method. The state
solve and the adjoint construction are implemented in the library; we study only the
\emph{compute cost}, \emph{memory}, and \emph{accuracy} of the gradient, and how
each scales with the number of parameters $N$. We contrast finite differences,
automatic differentiation of the loss (forward and reverse), the hand-written
adjoint, and AD equipped with a rule with different AD libraries in Julia.
\end{abstract}

\section{Methods}
\label{sec:methods}

The gradient maps the $N$ entries of the discretized potential $V$ to a scalar,
so the question that matters is how its cost grows with $N$.

\begin{itemize}
  \item \textbf{Finite differences (FD).} Perturb each entry in turn: $N{+}1$
  state solves, accuracy limited to $\sqrt{\varepsilon_{\mathrm{mach}}}$ by the
  step-size trade-off. Baseline.
  \item \textbf{Forward-mode AD.} One tangent direction per pass, so $N$ passes.
  It propagates a dual number through the Newton iteration rather than unrolling
  it to scalars, so a pass costs a constant times one state solve:
  machine-accurate, but $\mathcal O(N)$ like FD.
  \item \textbf{Reverse-mode AD.} A single backward pass returns the whole
  gradient, but it traces the Newton iteration and stores every intermediate
  state, so memory grows with the iteration count.
  \item \textbf{Adjoint (analytic).} One extra linear solve with the Newton-step
  operator at the converged solution: cost independent of $N$ and comparable to
  a single state solve. This is the reference.
  \item \textbf{Rule-based AD.} A custom rule for the state solve makes the AD
  tool compose the analytic adjoint instead of differentiating the iteration. It
  runs on the production solver and stays differentiable inside a larger pipeline.
\end{itemize}

\paragraph{Differentiating the linear solve.}
Every reverse-mode path carries a custom rule for the inner conjugate-gradient
solve, so the \emph{linear} system is never differentiated through; AD sees only
the \emph{Newton} iteration. (Forward mode simply propagates duals through the
fixed-iteration solver.) The naive--versus--ruled comparison then measures the
cost of unrolling Newton against the cost of replacing it by the adjoint rule,
not the easily avoided cost of unrolling CG.

\section{Theoretical complexity}
\label{sec:complexity}

Let $N$ be the number of parameters (grid points), $n_N$ the number of Newton
iterations, and $n_c$ the number of CG iterations per linear solve. One spectral
operator application costs $\mathcal O(N\log N)$, so one preconditioned CG solve
costs $T_{\mathrm{lin}}=\mathcal O(n_c\,N\log N)$ and one state (Newton) solve
costs $T_{\mathrm{sol}}=n_N\,T_{\mathrm{lin}}$. We express each gradient cost in
these units; \cref{tab:complexity} summarizes.

\begin{table}[t]
  \centering
  \caption{Cost of one gradient evaluation. $T_{\mathrm{sol}}=n_N\,T_{\mathrm{lin}}$
    is one state solve, $T_{\mathrm{lin}}=\mathcal O(n_c\,N\log N)$ one linear solve.}
  \label{tab:complexity}
  \begin{tabular}{lcccc}
    \toprule
    Method & State solves & Compute & Memory & Accuracy \\
    \midrule
    Finite differences          & $N{+}1$      & $\mathcal O(N\,T_{\mathrm{sol}})$ & $\mathcal O(N)$      & $\mathcal O(\sqrt{\varepsilon_{\mathrm{mach}}})$ \\
    Forward AD (duals)           & $\sim N$     & $\mathcal O(N\,T_{\mathrm{sol}})$ & $\mathcal O(N)$      & machine, at iterate \\
    Reverse AD naive (CG-ruled)  & $1 + $ bwd   & $\mathcal O(T_{\mathrm{sol}})$    & $\mathcal O(n_N N)$  & machine, at iterate \\
    Adjoint                      & $1$          & $\mathcal O(T_{\mathrm{sol}})$    & $\mathcal O(N)$      & exact, at $u^\star$ \\
    Rule-based AD                & $1$          & $\mathcal O(T_{\mathrm{sol}})$    & $\mathcal O(N)$      & $=$ adjoint \\
    \bottomrule
  \end{tabular}
\end{table}

\paragraph{Two regimes.}
FD and forward-mode AD cost $\mathcal O(N)$ state solves --- the price of probing
one input direction at a time. Reverse-mode AD, the adjoint and rule-based AD
return the \emph{entire} gradient at $\mathcal O(1)$ state solves, comparable to a
single loss evaluation and independent of $N$: the classical reverse-mode
result~\cite{johnson,levitt}.

\paragraph{Within the cheap group.}
What separates the $\mathcal O(1)$-solve methods is the constant and the memory.
The hand adjoint adds one linear solve and stores nothing extra. Naive
reverse-mode redoes the iteration backward --- $\approx n_N$ extra linear solves
($\approx 2\times$ the adjoint) and a tape of the $n_N$ Newton iterates
($\mathcal O(n_N N)$ memory). Rule-based AD injects exactly that adjoint,
recovering the single-solve cost and $\mathcal O(N)$ memory while staying
composable. On accuracy, forward- and reverse-mode AD differentiate the program
--- the $n_N$-th iterate $u_{n_N}(V)$ --- returning the exact derivative of an
approximate solution, off by $\mathcal O(\text{Newton residual})$; the adjoint
and the state-solve rule differentiate the converged solution through one exact
linear solve, matching to the linear-solve tolerance.

\section{Protocol}
\label{sec:protocol}

We fix a well-conditioned instance --- a forcing with nonzero mean, so the state
stays bounded away from zero --- and a single evaluation point $V_0$. Before any
timing, every gradient is checked against finite differences, so a method that is
fast only because it is wrong is excluded. \emph{Time} is the minimum over
warmed-up samples (BenchmarkTools, JIT excluded); \emph{memory} is the allocation
per call; \emph{accuracy} is the relative $\ell_2$ deviation from the adjoint
gradient. Times are also read in units of one state solve, the natural unit of
\cref{tab:complexity}. The study runs over $N=5,\dots,500$. The AD tools are
ForwardDiff (forward mode) and Mooncake and Zygote (reverse mode); the
reverse-mode paths use the CG rule, and the rule-based paths add the state-solve
rule.

\section{Results}
\label{sec:results}

\paragraph{Accuracy.}
\Cref{tab:accuracy,fig:accuracy} report the relative $\ell_2$ deviation of each
gradient from the analytic adjoint reference at $N=500$.

\begin{table}[t]
  \centering
  \caption{Relative $\ell_2$ deviation of each gradient from the analytic
    adjoint reference.}
  \label{tab:accuracy}
  \inctbl{gradient_data/accuracy_table.tex}
\end{table}

\begin{figure}[t]
  \centering
  \incfig{gradient_data/accuracy_vs_adjoint.pdf}
  \caption{Relative $\ell_2$ error of each method against the analytic adjoint (log scale).}
  \label{fig:accuracy}
\end{figure}

\noindent\fillin{State the FD deviation, the naive-AD deviation (the iterate gap,
of order the Newton residual), and confirm that rule-based AD sits at the adjoint
level rather than the iterate level.}

\paragraph{Cost at fixed size.}
\Cref{tab:cost} gives wall time and allocations at $N=500$, with each time
also as a multiple of one state solve.

\begin{table}[t]
  \centering
  \caption{Time and allocations per gradient evaluation at fixed $N$.}
  \label{tab:cost}
  \inctbl{gradient_data/cost_table.tex}
\end{table}

\noindent\fillin{Quote the adjoint and rule-based times in units of one solve
($\approx 1$--$2$), the FD and forward-AD slowdown ($\sim N$ solves), and the
reverse-naive allocation relative to the adjoint, where the $\mathcal O(n_N N)$
tape shows up even when the time is competitive.}

\paragraph{Scaling in $N$.}
\Cref{fig:scaling-time,fig:scaling-mem} show time and memory against $N$ on
log--log axes; the per-$N$ tables below give the backing data, one table per
grid size, each reporting the relative $\ell_2$ error against the analytic adjoint.

\begin{figure}[t]
  \centering
  \begin{minipage}{0.49\linewidth}\centering
    \incfig[\linewidth]{gradient_data/scaling_time.pdf}
    \caption{Time versus $N$.}\label{fig:scaling-time}
  \end{minipage}\hfill
  \begin{minipage}{0.49\linewidth}\centering
    \incfig[\linewidth]{gradient_data/scaling_memory.pdf}
    \caption{Allocations versus $N$.}\label{fig:scaling-mem}
  \end{minipage}
\end{figure}

\inctbl{gradient_data/scaling_table.tex}

\noindent\fillin{Read off the slopes: FD and forward-AD have unit slope in time
($\mathcal O(N)$ solves); reverse-mode, adjoint and rule-based are flat. In
memory, reverse-naive grows with $N$ while the others stay flat. Give the
crossover $N$ beyond which the adjoint is decisively cheaper than FD.}

\section{Discussion}
\label{sec:discussion}

The measurements confirm \cref{tab:complexity}. Finite differences and
forward-mode AD scale as $\mathcal O(N)$ state solves and are usable only as
baselines. Reverse-mode AD, the adjoint, and rule-based AD all compute the full
gradient at the cost of $\mathcal O(1)$ solves; within that group the hand
adjoint is the cheapest and lightest, naive reverse-mode pays the
Newton-unrolling factor in both compute ($\approx 2\times$) and memory
($\mathcal O(n_N N)$), and rule-based AD recovers the adjoint's efficiency while
remaining composable. For production the hand adjoint is the clear choice;
rule-based AD is preferable when the gradient must chain with other
differentiated code; naive differentiation is kept only as a baseline and a
correctness cross-check.

\paragraph{Reproducibility.}
\IfFileExists{gradient_data/meta.tex}{13th Gen Intel(R) Core(TM) i7-1360P (16 threads); Julia 1.12.6, BLAS threads 8, ForwardDiff 1.3.3, Mooncake 0.5.29, Zygote 0.7.10; commit b3ff9fc, 2026-06-16.}{\fillin{run gradient.jl to stamp version/hardware provenance here}}

\begin{thebibliography}{9}
  \bibitem{johnson} S.~G.~Johnson, \emph{Notes on Adjoint Methods for 18.335},
    MIT, 2006 (updated 2021).
  \bibitem{levitt} A.~Levitt, \emph{The Adjoint Trick}, 2020.
\end{thebibliography}

\end{document}